\documentclass[a4paper,12pt]{article}

\usepackage{graphicx}
\usepackage{hyperref}
\usepackage[italian]{babel}
\usepackage{parskip}
\usepackage{multicol}
\usepackage{verbatim}
\usepackage[a4paper,left=2cm,right=2cm,top=2.5cm,bottom=2.5cm]{geometry}
\usepackage{float}
\usepackage{setspace}

\setlength{\parskip}{1em}
\graphicspath{ {./images/} }
\hypersetup {
    colorlinks=true,
    linkcolor=black,
    pdfborderstyle={/S/U/W 1},
    }

    \begin{document}
    \begin{titlepage}
        \begin{center}
            \begin{minipage}{0.3\textwidth}
                \includegraphics[width=\linewidth]{Logo_Universita_Padova300x300.webp}
            \end{minipage}

        \end{center}

        \begin{center}
            \noindent\rule{0.7\textwidth}{0.4pt}\\[0.5cm]
            \large\textbf{Università degli Studi di Padova}\\
            \large Corso di Laurea in Informatica\\
            \small A. A. 2025/26
        \end{center}

        \vspace{1,5cm}    

        \begin{center}
            \begin{minipage}{0.5\textwidth}
                \includegraphics[width=\linewidth]{logo.webp}
            \end{minipage}
        \end{center}

        \vspace{1,5cm}    
        \begin{center}
            {\Large \textbf{Relazione}\par}
            \vspace{0.6cm}
            {\Large Progetto di Tecnologie Web\par}
            \vspace{0.3cm}
        \end{center}
        \vfill

        \centering
        \begin{tabular}{r|l}
            \textbf{Studenti} & Sofia De Blasi - matricola - email \\
            & Nenad Radulovic - matricola - email \\
            & Giacomo Speggiorin - matricola - email \\
            & Nicola Simionato - matricola - email \\
            & \\
            \textbf{Indirizzo del sito} & indirizzo\\
            & \\
            \textbf{Utente} & 
            \textbf{Username:} user, \textbf{Password:} user \\ 


        \end{tabular}

        \vfill

    \end{titlepage}     
    \newpage
    \tableofcontents
    \newpage

\section{Introduzione}
    \subsection{Com'è nata l'idea}
    L'idea di questo progetto nasce da una domanda che ci siamo posti più volte nella nostra quotidianità: \textit{``Cosa potremmo fare questo weekend?''}.  
    Nelle città universitarie vengono spesso organizzati eventi e attività dedicate agli studenti, ma è seccante doverli andare a cercare ogni volta. Bisogna innanzitutto sapere cosa si sta cercando (il che non è affatto scontato) e, successivamente, sfogliare di volta in volta le pagine social dei vari organizzatori, che sono sempre diversi e non centralizzati.
    
    Da qui è nata la riflessione su quanto sarebbe comodo se fossero gli organizzatori stessi ad andare incontro agli studenti, e su quanto si semplificherebbe la ricerca se esistesse un'unica bacheca dove tutti potessero ``appendere'' i propri volantini.
    
    Mentre sviluppavamo l’idea ci siamo resi conto che il problema non riguarda solo gli eventi: esistono molti altri servizi di cui gli studenti hanno bisogno, ma che non hanno un punto di riferimento chiaro su dove trovarli. Da questa consapevolezza è nato il desiderio di creare uno spazio unico, accessibile e organizzato, pensato per raccogliere tutto ciò che può essere utile nella vita quotidiana universitaria.

    \subsection{L'obiettivo di StudentSpace}
    StudentSpace propone una \textit{bacheca online} che raccolga in un unico luogo eventi, attività e servizi di interesse per gli studenti, andando oltre il contesto strettamente accademico.  
    Attraverso la piattaforma sarà possibile trovare eventi sportivi, attività serali, concerti e altre occasioni di socialità. Gli studenti potranno inoltre offrire la propria disponibilità per ripetizioni, oppure consultare annunci di affitti relativi alla propria città.
    
    L'obiettivo principale è garantire a studenti e studentesse un accesso rapido, centralizzato e completo a ciò che potrebbe interessarli nella vita quotidiana universitaria.  
    Allo stesso tempo, la piattaforma intende agevolare gli organizzatori che si rivolgono a questo target, offrendo loro un canale unico e mirato per raggiungere un pubblico ampio, omogeneo e specifico.

\section{Analisi del lavoro svolto}
    \subsection{Processo decisionale}
    Durante tutto lo sviluppo del progetto ci si è trovati a dover prendere diverse decisioni in merito a design, funzionalità, contenuto ecc. 
    Nella maggior parte dei casi questi punti sono stati discussi in presenza, tra tutti i membri del gruppo. 
    Ovviamente, questo metodo comporta delle chiare limitazioni, si è quindi deciso di adottare
    anche dei metodi di comunicazione (sia sincroni che asincroni) alternativi.\\\\
    Per quanto riguarda le comunicazioni asincrone, è stata usata ampiamente la 
    piattaforma di messaggistica \textbf{\underline{\hyperref[sec:strumenti]{Whatsapp}}}. 
    Inoltre, si è deciso di adottare \textbf{\underline{\hyperref[sec:strumenti]{Discord}}} come strumento di comunicazione asincrona, in supporto alle riunioni in presenza.

    \subsection{Modalità di lavoro}
    Per adattarsi meglio alle esigenze del gruppo si è deciso che i membri del gruppo avrebbero svolto ciascuno le proprie attività in maniera prevalentemente asincrona.
    Questa decisione è stata presa per riuscire a dare più libertà ai singoli membri nell'organizzare il proprio lavoro. 
    La scelta è stata inoltre supportata dalla volontà del gruppo di dare maggiore fluidità allo sviluppo del progetto. 
    Il lavoro è stato infatti suddiviso in modo che ogni membro andasse a lavorare su punti diversi rispetto agli altri, cosicché nessuno dovesse preoccuparsi che il proprio lavoro andasse in conflitto con quello di qualcun altro. 
    A sostegno di ciò si è deciso di implementare un sistema di versionamento del codice (con annesso hosting online) tramite \textbf{\underline{\hyperref[sec:strumenti]{Git}}} e \textbf{\underline{\hyperref[sec:strumenti]{Github}}}.

    \subsection{Divisione del lavoro}
    Visto lo scopo di natura accademica del progetto si è ritenuto opportuno che
    ciascun membro del gruppo fosse coinvolto (quantomeno in parte) in ogni ambito del progetto.
    Questo ha permesso a ogni membro di avere un idea generale sul funzionamento del 
    programma, e di mettere mano direttamente a ogni fase del suo sviluppo.\\
    Periodicamente sono state svolte delle riunioni di assestamento, nelle quali veniva
    aggiornata la lista delle cose da fare per il periodo sucessivo e alla fine di ogni riunione le attività da svolgere venivano assegnate ai membri in modo equilibrato.
    Per facilitare questo processo si è fatto largo uso delle Github Issue.

    \section{Progettazione}
\subsection{Tipologie di utenti}

Il sistema prevede due tipologie principali di utenti, differenziate per permessi e funzionalità disponibili.

\begin{itemize}

    \item \textbf{Utente non registrato:} può esplorare liberamente le categorie e leggere gli annunci.
    \begin{itemize}
        \item \textbf{Azioni e permessi:} accesso completo alle pagine \textit{index} ed \textit{esplora}, con possibilità di effettuare ricerche e applicare filtri; apertura del dettaglio dell'\textit{annuncio} e della pagina di \textit{profilo pubblico} del proponente; possibilità di avviare la procedura di registrazione o accesso.
        \item \textbf{Limitazioni:} non può pubblicare, gestire o salvare annunci nei preferiti. Se tenta di eseguire azioni riservate, il sistema lo reindirizza alla pagina di accesso, offrendo anche un collegamento diretto alla pagina di registrazione qualora non possedesse un account.
    \end{itemize}

    \item \textbf{Utente registrato:} oltre alle funzionalità dell’utente non registrato, può pubblicare e gestire i propri annunci, salvare annunci altrui nei preferiti e accedere alla propria area personale.
    \begin{itemize}
        \item \textbf{Azioni e permessi:} creazione, modifica ed eliminazione dei propri annunci; gestione del profilo personale; consultazione degli annunci preferiti e degli annunci pubblicati.
        \item \textbf{Limitazioni:} non può modificare annunci altrui; è soggetto a vincoli relativi ai campi obbligatori e ai formati previsti, con validazione sia lato client che lato server.
    \end{itemize}

\end{itemize}


    \subsection{Analisi dei contenuti} % Tipi di annunci, Campi generici, Campi specifici per categoria
    Il sito web si propone come una bacheca universitaria online, in cui gli utenti possono
    pubblicare e consultare annunci suddivisi in quattro categorie principali: \textbf{Affitti},
    \textbf{Esperimenti}, \textbf{Eventi} e \textbf{Ripetizioni}. 
    L’analisi dei contenuti si concentra sui tipi di annunci, sui campi comuni 
    a tutte le categorie e sui campi specifici per ciascuna categoria.

    \subsubsection{Tipi di annunci}
    \begin{itemize}
        \item \textbf{Affitti}: annunci relativi a stanze, appartamenti o posti letto disponibili per studenti.
        \item \textbf{Esperimenti}: annunci riguardanti attività di laboratorio, progetti di ricerca o collaborazioni scientifiche.
        \item \textbf{Eventi}: annunci di conferenze, seminari, feste universitarie, attività culturali o sportive.
        \item \textbf{Ripetizioni}: annunci di studenti che offrono supporto didattico su materie specifiche.
    \end{itemize}

    \subsubsection{Campi generici}
    Ogni annuncio, indipendentemente dalla categoria, contiene i seguenti campi:
    \begin{itemize}
        \item \textbf{Categoria}: selezionata dall’utente tra le quattro disponibili.
        \item \textbf{Titolo}: breve descrizione sintetica dell’annuncio.
        \item \textbf{Data di pubblicazione}: generata automaticamente dal sistema.
        \item \textbf{Città}: riferimento all’utente registrato che ha pubblicato l’annuncio.
        \item \textbf{Descrizione}: testo libero con i dettagli dell’annuncio. 
        \item \textbf{Autore}: riferimento all’utente registrato che ha pubblicato l’annuncio. 
    \end{itemize}

    \subsubsection{Campi specifici per categoria}
    Oltre ai campi generici, ciascuna categoria prevede attributi specifici (3 per categoria):
    \begin{multicols}{2}
        \begin{itemize}
            \item \textbf{Affitti}:
                \begin{itemize} % decidere se tenere tutti questi attrubuti
                    \item Prezzo mensile
                    \item Indirizzo
                    \item Numero coinquilini
                \end{itemize}
            \item \textbf{Esperimenti}:
                \begin{itemize}
                    \item Laboratorio di riferimento
                    \item Durata prevista
                    \item Compenso
                \end{itemize}
            \item \textbf{Eventi}:
                \begin{itemize}
                    \item Data Eventi
                    \item Luogo
                    \item Costo entrata
                \end{itemize}
            \item \textbf{Ripetizioni}:
                \begin{itemize}
                    \item Materia
                    \item Livello (scuola superiore, università, ecc.)
                    \item Prezzo orario
                \end{itemize}
        \end{itemize}
    \end{multicols}
    \subsection{Descrizione delle funzionalità} % registrazione/login, pubblicazione annunci, modifica/eliminazione, filtraggio annunci, validazione input (client/server), accessibilità e responsività
    Il sistema offre una serie di funzionalità fondamentali per soddisfare i requisiti 
    del progetto e le esigenze degli utenti.

    \subsubsection{Registrazione e login}
    \begin{itemize}
        \item Registrazione di nuovi utenti: è permessa, per gli utenti che volessero ampliare la loro esperienza di utilizzo, la creazione di un account. La registrazione richiederà all'utente di fornire le seguenti informazioni:
            \begin{itemize}
                \item Nome
                \item Cognome
                \item Email
                \item Citta: intesa come città universitaria di riferimento.
                \item Password
            \end{itemize}
        \item Login: processo di identificazione degli utenti registrati che sblocca tutte le funzionalità.
        \item Qualora l'utente inserisca dei dati non validi verrà bloccato l'invio della
            form e verranno mostrati dei messaggi di errore (si veda \underline{\hyperref[sviluppo:validazione]{la parte relativa alla validazione}}).

    \end{itemize}

    \subsubsection{Pubblicazione annunci}
    \begin{itemize}
        \item È possibile, per gli utenti registrati, pubblicare annunci compilando un'apposito form. 
            Sono necessari per la pubblicazione:
            \begin{itemize}
                \item Campi generici: titolo, descrizione e categoria.
                \item Campi specifici mostrati dinamicamente in base alla categoria scelta.
            \end{itemize}
        \item In seguito alla pubblicazione sarà possibile risalire all'autore dell'annuncio.
        \item Qualora l'utente inserisca dei dati non validi verrà bloccato l'invio della
            form e verranno mostrati dei messaggi di errore (si veda \underline{\hyperref[sviluppo:validazione]{la parte relativa alla validazione}}).

    \end{itemize}

    \subsubsection{Modifica, eliminazione e salvataggio degli annunci}  
    \begin{itemize}
        \item Un utente registrato può decidere di salvare degli annunci aggiungendoli 
            alla propria lista dei preferiti.
        \item L'autore di un annuncio può decidere di modificare o cancellare i propri annunci.
        \item Un annuncio non apparirà più all'interno della lista dei preferiti degli utenti 
            dopo essere stato eliminato dal sistema.
    \end{itemize}

    \subsubsection{Filtraggio e ricerca degli annunci}

\begin{itemize}

    \item Il sistema permette di applicare diversi filtri per affinare la ricerca degli annunci. 
    Sono previsti:
    \begin{itemize}
        \item \textbf{Filtri generici}: categoria e data di pubblicazione.
        \item \textbf{Filtri specifici per categoria}: campi aggiuntivi che variano in base alla tipologia di annuncio.
    \end{itemize}

    \item È inoltre disponibile una ricerca testuale che opera sul contenuto del titolo, della descrizione e sull'email dell'autore dell'annuncio.

    \item Qualora vengano inseriti nei campi liberi valori non validi o potenzialmente compromettenti, il sistema li ignora senza generare errori, garantendo comunque un'esperienza di ricerca stabile e coerente.
    
\end{itemize}


    \subsection{Archietettura del sistema} % struttura delle cartelle, separazione contenuto/presentazione/comportamento, panoramica componenti

    \subsection{Struttura del sito} %mappa delle pagine, funzioni, URL
    \subsection{Mockup}
    \subsection{Responsive Web Design e Mobile First} 

    \newpage
    \section{Sviluppo} % decidere struttura
    \begin{comment}
        \subsection{HTML} % sementica
        \subsection{CSS}
        \subsection{JavaScript} % validazione client + interazione
        \subsection{PHP}
        \subsection{Database}
        \subsection{Errori}
        \subsection{Validazione} 
        \subsection{Gestione delle immagini} 
        \subsection{Palette colori e font} 
        \subsection{Configurazione .htaccess} % password hashing, SQL injection, sanitizzazione input, protezione aree riservate, gestione sessioni
    \end{comment}
    % o struttura più sintetica

    \subsection{Database} 
    \subsubsection{Analisi dei requisiti del database e progettazione logica}

    La base di dati deve supportare la gestione degli utenti e degli annunci pubblicati
    sulla bacheca universitaria. In particolare, deve garantire la memorizzazione dei
    campi generici comuni a tutte le categorie di annunci e dei campi specifici per ciascuna
    categoria. Inoltre, deve consentire operazioni di
    inserimento. %, modifica e cancellazione.

    \paragraph{Entità principali}
    \begin{itemize}
        \item \textbf{Utente}: rappresenta gli utenti registrati alla piattaforma.\\
            \texttt{IdUtente(PK)}, \texttt{Nome}, \texttt{Cognome}, \texttt{Email (UNIQUE)}, \texttt{Password}, \texttt{IdCitta (FK)*}.

            \begin{comment} 
            \item \textbf{Categoria}: rappresenta le tipologie di annuncio.\\
                \texttt{IdCategoria(PK)}, \texttt{NomeCategoria (UNIQUE)}.
            \end{comment}

        \item \textbf{Città}: rappresenta la sede universitaria o città di riferimento  di un utente o di un annuncio.\\
            \texttt{IdCitta(PK)}, \texttt{NomeCitta (UNIQUE)}.

        \item \textbf{Annuncio}: rappresenta un contenuto pubblicato sulla bacheca.\\
            \texttt{IdAnnuncio(PK)}, \texttt{Titolo}, \texttt{Descrizione}, \texttt{DataPubblicazione (DEFAULT CURRENT\_TIMESTAMP)}, \texttt{IdUtente (FK)}, \texttt{NomeCategoria}, \texttt{IdCitta (FK)}.

        \item \textbf{Specializzazioni di Annuncio}: tabelle figlie che estendono Annuncio con 3 attributi specifici.
            \begin{itemize}
                \item \textbf{Affitti}: \texttt{IdAnnuncio (PK, FK)}, \texttt{PrezzoMensile}, \texttt{Indirizzo}, \texttt{NumeroInquilini}.
                \item \textbf{Esperimenti}: \texttt{IdAnnuncio (PK, FK)}, \texttt{Laboratorio}, \texttt{DurataPrevista}, \texttt{Compenso}.
                \item \textbf{Event}: \texttt{IdAnnuncio (PK, FK)}, \texttt{DataEventi}, \texttt{Luogo}, \texttt{CostoEntrata}.
                \item \textbf{Ripetizioni}: \texttt{IdAnnuncio (PK, FK)}, \texttt{Materia}, \texttt{Livello}, \texttt{PrezzoOrario}.
            \end{itemize}

        \item \textbf{ImmaginiAnnuncio}: rappresenta le immagini associate ad un annuncio.\\
            \texttt{IdImmagine(PK)}, \texttt{IdAnnuncio (FK)}, \texttt{Percorso}, \texttt{AltText*}, \texttt{Decorativa}, \texttt{Ordine}.

    \end{itemize}

    \paragraph{Relazioni tra entità}

    \begin{itemize}
        \item \textbf{Pubblicazione:} collega \texttt{Utente} e \texttt{Annuncio}.\\
            \texttt{Utente (1,1)} $\longrightarrow$ \texttt{Annuncio (0,N)}\\
            Ogni annuncio è pubblicato da un solo utente, mentre un utente può pubblicare più annunci.

            \begin{comment}
            \item \textbf{Appartenenza:} collega \texttt{Categoria} e \texttt{Annuncio}.\\
                \texttt{Categoria (1,1)} $\longrightarrow$ \texttt{Annuncio (0,N)}\\
                Ogni annuncio appartiene ad una sola categoria, mentre una categoria può essere associata a più annunci.
            \end{comment}

        \item \textbf{Collocazione:} collega \texttt{Città} e \texttt{Annuncio}.\\
            \texttt{Città (1,1)} $\longrightarrow$ \texttt{Annuncio (0,N)}\\
            Ogni annuncio è associato ad una sola città, mentre una città può avere più annunci.

        \item \textbf{SedeUniversitaria:} collega \texttt{Città} e \texttt{Utente}.\\
            \texttt{Città (1,1)} $\longrightarrow$ \texttt{Utente (0,N)}\\
            Ogni utente è associato ad una sola città, mentre una città può avere più utenti.

        \item \textbf{Associazione immagini:} collega \texttt{Annuncio} e \texttt{ImmaginiAnnuncio}.\\
            \texttt{Annuncio (1,1)} $\longrightarrow$ \texttt{ImmaginiAnnuncio (0,N)}\\
            Ogni annuncio può avere da zero a più immagini, mentre ogni immagine è riferita ad un solo annuncio.

    \end{itemize}


    \subsubsection{Schema ER}
    \begin{figure}[H] 
        \centering
        \includegraphics[width=1\textwidth]{ERCatAttr.jpg}
        \caption{Schema ER}
        \label{fig:ER}
    \end{figure}

    \paragraph{Generalizzazione IS-A}
    Nel modello concettuale iniziale, l’entità \texttt{Annuncio} fungeva da padre nella generalizzazione per le quattro categorie specifiche figlie (\texttt{Affitti}, \texttt{Esperimenti}, \texttt{Eventi}, \texttt{Ripetizioni}).  
    In fase di progettazione logica si è ristrutturato lo schema ER con una \textbf{sostituzione della generalizzazione con associazioni}.
    \\ Questa scelta consente di: evitare la ripetizione dei campi generici, mantenere la normalizzazione, facilitare l’estendibilità del modello, separare logicamente i dati specifici per categoria. \\ Tuttavia aggiunge un vincolo, ossia: ogni occorrenza del padre \texttt{Annuncio} non può partecipare contemporaneamente a due relazioni \texttt{IS-[categoria]}. 

    \subsubsection{Schema ER ristrutturato}
    \begin{figure}[H] 
        \centering
        \includegraphics[width=1\textwidth]{ERristruttCatAttr.jpg}
        \caption{Schema ER ristrutturato}
        \label{fig:ER}
    \end{figure}

    \textbf{Relaazione IS-[categoria]:} collega \texttt{Annuncio} alle specializzazioni (\texttt{Affitti}, \texttt{Esperimenti}, \texttt{Eventi}, \texttt{Ripetizioni}).\\
    \texttt{Annuncio (1,1)} $\longrightarrow$ \texttt{Specializzazione (0,1)}\\
    Ogni annuncio può avere al massimo una specializzazione, e ogni specializzazione è riferita ad un solo annuncio.

    \subsubsection{Scelta delle chiavi primarie}
    Nella progettazione del database si è deciso di utilizzare come chiavi primarie codici (ID numerici auto\-incrementali). 

    Questa scelta è motivata da diversi fattori:
    \begin{itemize}
        \item \textbf{Stabilità}: gli attributi naturali (es. \texttt{Email}, \begin{comment} \texttt{NomeCategoria} \end{comment}, \texttt{NomeCitta}) possono cambiare nel tempo o essere modificati dall’utente, mentre un ID rimane invariabile.
            \item \textbf{Efficienza}: gli ID numerici sono più leggeri da gestire negli indici e nelle join rispetto a stringhe lunghe come email o nomi.
            \item \textbf{Semplicità}: l’uso di codici numerici rende più chiaro e uniforme lo schema, evitando di dover scegliere chiavi composite o troppo complesse.
            \item \textbf{Integrità}: gli attributi naturali vengono comunque vincolati con \texttt{UNIQUE} (es. \texttt{Email}, \begin{comment} \texttt{NomeCategoria} \end{comment}, \texttt{NomeCitta}) per garantire l’unicità, ma non sono usati come chiavi primarie.
                \item \textbf{Estendibilità}: se in futuro si aggiungessero nuove categorie o sedi, l’uso di ID permetterebbe di gestire facilmente l’espansione senza modificare le relazioni esistenti.
    \end{itemize}

    \subsubsection{Schema logico (relazionale)}
    Di seguito è riportato lo schema logico, in cui l'asterisco dopo il nome di un attributo indica che tale attributo può assumere valori nulli.

    \begin{itemize}
        \item \textbf{Utente}(\underline{IdUtente}, Nome, Cognome, Email [UNIQUE], Password, SedeUniversitaria*
            \begin{itemize}
                \item Utente.SedeUniversitaria $\rightarrow$ Citta.IdCitta
            \end{itemize}

        \begin{comment} \item \textbf{Categoria}(\underline{IdCategoria}, NomeCategoria [UNIQUE]) \end{comment}

            \item \textbf{Citta}(\underline{IdCitta}, NomeCitta [UNIQUE])

            \item \textbf{Annuncio}(\underline{IdAnnuncio}, Titolo, Descrizione, DataPubblicazione, Autore, Categoria, \\ Collocazione)
                \begin{itemize}
                    \item Annuncio.Autore $\rightarrow$ Utente.IdUtente
                    \begin{comment} \item Annuncio.Categoria $\rightarrow$ Categoria.IdCategoria \end{comment}
                        \item Annuncio.Collocazione$\rightarrow$ Citta.IdCitta
                \end{itemize}

            \item \textbf{Affitti}(\underline{IdAnnuncio}, PrezzoMensileAff, IndirizzoAffitti, NumeroCoinquilini)
                \begin{itemize}
                    \item Affitti.IdAnnuncio$\rightarrow$ Annuncio.IdAnnuncio
                \end{itemize}

            \item \textbf{Esperimenti}(\underline{IdAnnuncio}, LaboratorioEsp, DurataPrevEsp, CompensoEsp)
                \begin{itemize}
                    \item Esperimenti.IdAnnuncio$\rightarrow$ Annuncio.IdAnnuncio
                \end{itemize}

            \item \textbf{Eventi}(\underline{IdAnnuncio}, DataEventi, LuogoEventi, CostoEntrataEventi)
                \begin{itemize}
                    \item Eventi.IdAnnuncio$\rightarrow$ Annuncio.IdAnnuncio
                \end{itemize}

            \item \textbf{Ripetizioni}(\underline{IdAnnuncio}, MateriaRipetizioni, LivelloRip, PrezzoOrarioRip)
                \begin{itemize}
                    \item Ripetizioni.IdAnnuncio$\rightarrow$ Annuncio.IdAnnuncio
                \end{itemize}

            \item \textbf{ImmaginiAnnuncio}(\underline{IdImmagine}, IdAnnuncio (FK), Percorso, AltText*, Decorativa, Ordine)
                \begin{itemize}
                    \item ImmaginiAnnuncio.IdAnnuncio $\rightarrow$ Annuncio.IdAnnuncio
                \end{itemize}

    \end{itemize}

    \subsubsection{Normalizzazione}
    Lo schema progettato rispetta la terza forma normale (3NF):
    \begin{itemize}
        \item \textbf{1NF}: tutti gli attributi sono atomici e non esistono campi multivalore o ripetuti.

        \item \textbf{2NF}: ogni tabella ha una chiave primaria semplice (IdUtente, \begin{comment} IdCategoria \end{comment}, IdCitta, IdAnnuncio). Nelle specializzazioni (\texttt{Affitti}, \texttt{Esperimenti}, \texttt{Eventi}, \texttt{Ripetizioni}) la chiave primaria è \texttt{IdAnnuncio} e tutti gli attributi dipendono interamente da essa.

            \item \textbf{3NF}: non esistono dipendenze transitive tra attributi non chiave.
    \end{itemize}

    \subsubsection{Vincoli di integrità}

    Lo schema logico rispetta i principali vincoli di integrità:

    \begin{itemize}
        \item \textbf{Unicità}: Email utente, \begin{comment} NomeCategoria \end{comment}, NomeCitta.
            \item \textbf{Dominio}: campi obbligatori definiti come NOT NULL.
            \item \textbf{Referenziali}: ON DELETE CASCADE sulle specializzazioni; SET NULL in Annuncio su Utente, Categoria, Citta nel caso vengano eliminati.
            \item \textbf{Generalizzazione IS-A}: ogni annuncio può avere al massimo una specializzazione coerente con la categoria.
    \end{itemize}



    \subsection{Front-end (HTML, CSS, JS)}
    \subsubsection{HTML}
    Nel documento è stato inserito il tag:
    \begin{verbatim}
    <meta name="viewport" content="width=device-width"/>
    \end{verbatim}
    Tale configurazione permette ai fogli di stile di riferirsi alla dimensione effettiva del dispositivo e non alla sua risoluzione fisica. In questo modo si evita che, su smartphone ad alta densità di pixel, venga caricata la visualizzazione desktop, garantendo un comportamento responsivo coerente con il paradigma \textit{mobile first}.

    Nella pagina \textit{esplora} è stato scelto di utilizzare l’elemento semantico \texttt{<aside>} per contenere i filtri di ricerca, invece di un elemento \texttt{<section>}. Tale scelta deriva dal fatto che i filtri rappresentano contenuti accessori rispetto al flusso principale della pagina, costituito dall’elenco degli annunci. L’uso di \texttt{<aside>} migliora la chiarezza semantica e supporta una migliore interpretazione da parte delle tecnologie assistive.

    Per quanto riguarda le singole card, gestite tramite \texttt{cardTemplate.html}, è stato deciso di utilizzare elementi \texttt{<h3>} per i titoli degli annunci, anziché una struttura basata su \texttt{<dt>} e \texttt{<dd>}, come nel seguente esempio:
    \begin{verbatim}
    <dt>Titolo:</dt>
    <dd>[TitoloAnnuncio]</dd>
    \end{verbatim}
    L’uso degli heading consente di riflettere correttamente la gerarchia logica dei contenuti: \texttt{<h2>} identifica la sezione principale (ad esempio l’elenco degli annunci), mentre \texttt{<h3>} rappresenta ciascun singolo annuncio. Questa scelta migliora la struttura semantica complessiva, favorisce un migliore posizionamento nei motori di ricerca e permette agli screen reader di interpretare più efficacemente l’organizzazione informativa della pagina.

    \subsubsection{CSS e decisioni di design (immagini, palette colori)} 


    \subsection{Back-end (PHP)}


    \subsection{Validazione e Sicurezza} \label{sviluppo:validazione}
    \begin{itemize}
        \item \textbf{Lato client}: controlli immediati con JavaScript (campi obbligatori, formati corretti).
        \item \textbf{Lato server}: verifica dati ricevuti prima dell’inserimento nel database.
        \item Messaggi di errore chiari e accessibili.
        \item non sono concessi tag negli input
    \end{itemize}


    \newpage
    \section{Test}
    \subsection{Test funzionali}
    Descrizione dei test su login, registrazione, pubblicazione e filtraggio annunci.

    \subsection{Test di compatibilità}
    Test su diversi browser (Chrome, Firefox, Edge, Safari) e dispositivi (desktop, tablet, smartphone).

    \subsection{Test di accessibilità}
    Color contrast, navigazione da tastiera, screen reader.

    \subsection{Requisiti tecnici}
    \subsubsection{Accessibilità e responsività}
    \begin{itemize}
        \item Layout realizzato con CSS3, Flexbox e Grid.\
        \item Separazione netta tra contenuto (HTML5), presentazione (CSS) e comportamento (JS/PHP).
        \item Compatibilità cross-browser e adattamento a diverse dimensioni di schermo.
        \item Uso di etichette semantiche, contrasto cromatico adeguato e font leggibili.

        \item Validazione lato client (formato email, lunghezza password) e lato server.
    \end{itemize}


    \section{Accessibilità}
    \subsection{Introduzione} 
    \subsection{WAI-ARIA e screen-reader} 
    \subsection{Progressive Enhancement} 
    \subsection{Supporto ai dispositivi “legacy”} 
    \subsection{Ulteriori funzionalità di accessibilità}

    \section{SEO e prestazioni} %meta tag, URL parlanti, sitemap, compressione immagini, lazy loading, minificazione
    \subsection{Analisi delle possibili ricerche su Google} % “Affitti stanza Padova”, “eventi Padova weekend”, “Ripetizioni informatica Padova”…)

    \newpage
    \section{Strumenti utilizzati} \label{sec:strumenti}
    \begin{comment}
        \setstretch{0.7}
        \subsection{Canva} \label{strumenti:canva}
        \underline{\url{https://www.canva.com/it_it/}} \\\\
        Piattaforma di progettazione grafica per creare contenuti visuali per presentazioni, siti web e prodotti simili. 

        \subsection{Discord} \label{strumenti:discord}
        \underline{\url{https://discord.com/}} \\\\
        Piattaforma di messaggistica istantanea, di distribuzione digitale e di voice over IP, utilizzata per la comunicazione all'interno server Discord.
        Gli utenti comunicano con chiamate vocali, videochiamate, messaggi di testo, media e file in chat private o come membri di un server Discord. Quest'ultimi sono una raccolta di canali di tipo vocale e/o testuale. 

        \subsection{Git} \label{strumenti:git}
        \underline{\url{https://git-scm.com/}} \\\\
        Applicazione software per il controllo di versione distribuito utilizzabile da interfaccia a riga di comando. Viene utilizzato il versionamento delle repository dedicate al codice e alla documentazione del progetto.

        \subsection{Github} \label{strumenti:github}
        \underline{\url{https://git-scm.com/}} \\\\
        È un servizio di hosting per progetti software che implementa lo strumento di controllo versione distribuito Git.  

        \subsection{LaTeX} \label{strumenti:latex}
        \underline{\url{https://www.latex-project.org/get/}} \\\\
        Linguaggio di marcatura del testo, usato ampiamente in ambito accademico per la preparazione di testi.

        \subsection{Silktide} \label{strumenti:silktide}
        \underline{\url{https://silktide.com/}} \\\\
        È un'estensione che facilita i test sull'accessibilità.

        \subsection{W3C Validator} \label{strumenti:w3cvalidator}
        \underline{\url{http://validator.w3.org/}} \\\\
        Validatore markup HTML/XHTML.

        \subsection{WhatsApp} \label{strumenti:whatsapp}
        \underline{\url{https://www.whatsapp.com/}} \\\\

        \subsection{WebAIM} \label{strumenti:webaim}
        \underline{\url{https://webaim.org/resources/contrastchecker/}} \\\\
    \end{comment}
    \begin{itemize}
        \item \textbf{Canva} 
            \underline{\url{https://www.canva.com/it_it/}} \\
            Piattaforma di progettazione grafica per creare contenuti visuali per presentazioni, siti web e prodotti simili.

        \item \textbf{Discord} 
            \underline{\url{https://discord.com/}} \\
            Piattaforma di messaggistica istantanea e distribuzione digitale, utilizzata per la comunicazione del gruppo. 
            Supporta chiamate vocali, videochiamate, messaggi di testo e condivisione file.

        \item \textbf{Git}
            \underline{\url{https://git-scm.com/}} \\
            Applicazione software per il controllo di versione distribuito, usata per il versionamento del codice e della documentazione.

        \item \textbf{Github}
            \underline{\url{https://github.com/}} \\
            Servizio di hosting per progetti software che implementa Git, utilizzato per la collaborazione e la condivisione del codice.

        \item \textbf{LaTeX}
            \underline{\url{https://www.latex-project.org/get/}} \\
            Linguaggio di marcatura del testo, usato per la preparazione della relazione accademica.

        \item \textbf{Silktide}
            \underline{\url{https://silktide.com/}} \\
            Estensione che facilita i test sull’accessibilità.

        \item \textbf{W3C Validator} 
            \underline{\url{http://validator.w3.org/}} \\
            Strumento per validare il markup HTML/XHTML.

        \item \textbf{WhatsApp} 
            \underline{\url{https://www.whatsapp.com/}} \\
            Piattaforma di messaggistica istantanea usata per la comunicazione rapida tra i membri del gruppo.

        \item \textbf{WebAIM}  
            \underline{\url{https://webaim.org/resources/contrastchecker/}} \\
            Strumento online per verificare il contrasto dei colori e migliorare l’accessibilità.
    \end{itemize}


    \section{Guida all’installazione del progetto}


\end{document}
